\documentclass[11pt]{article}
\usepackage[margin=1in]{geometry}
\usepackage{booktabs,tabularx,longtable,array}
\usepackage{hyperref}
\usepackage{enumitem}
\title{Supplementary Information\\Cost--Resilience Thresholds for Battery--Hydrogen Hybrid Microgrids in Islanded Power Systems}
\author{Jaewon Lee}
\date{}
\begin{document}
\maketitle

\section*{Supplementary file caption}
This supplementary file documents the fixed-portfolio definitions, synthetic-load construction, outage-setting rationale, screened parameter ranges, and reproducibility workflow used in the main manuscript.

\section*{S1. Fixed portfolio definitions and design intent}
Table~S1 reports the fixed capacities used in the three candidate portfolios. These values are taken directly from the scenario configuration used to generate the submitted results.

\begin{longtable}{p{2.2cm}p{2.5cm}rrrrrr}
\caption{Fixed capacities used in the candidate portfolios.}\\
\toprule
Portfolio & Label & PV (MW) & Battery power (MW) & Battery energy (MWh) & Diesel (MW) & Electrolyzer (MW) & Fuel cell (MW) \\
\midrule
\endfirsthead
\toprule
Portfolio & Label & PV (MW) & Battery power (MW) & Battery energy (MWh) & Diesel (MW) & Electrolyzer (MW) & Fuel cell (MW) \\
\midrule
\endhead
battery\_only & PV + battery & 7.0 & 2.2 & 24.0 & 0.0 & 0.0 & 0.0 \\
diesel\_battery & PV + diesel + battery & 2.8 & 1.5 & 8.0 & 2.4 & 0.0 & 0.0 \\
battery\_hydrogen & PV + battery + hydrogen & 14.0 & 2.0 & 12.0 & 0.0 & 5.0 & 2.2 \\
\bottomrule
\end{longtable}

The battery--hydrogen portfolio additionally includes a hydrogen tank with a fixed inventory capacity of 50{,}000 kg. The fixed-capacity design is intentional: it allows the manuscript to compare operational and resilience behavior across candidate portfolios under common assumptions without claiming endogenous optimal sizing.

\section*{S2. Synthetic load archetype and case-study inputs}
The case-study inputs were generated from the scenario file and the prepared 8760-hour dataset used by the dispatch model. The synthetic load archetype is calibrated using the following inputs:
\begin{itemize}[leftmargin=1.5em]
  \item peak demand = 1.65 MW
  \item load factor = 0.72
  \item evening peak strength = 0.30
  \item morning peak strength = 0.10
  \item commercial load fraction = 0.14
  \item tourism multiplier = 1.08 in March, April, and December
  \item seasonal monthly multipliers = [1.02, 1.03, 1.04, 1.05, 1.04, 1.02, 0.99, 0.98, 0.99, 1.00, 1.01, 1.03]
  \item random noise fraction = 0.025 with seed = 42
\end{itemize}

This construction is not intended to reproduce a licensed feeder's exact dispatch history. Its purpose is to provide a transparent island-like load profile that preserves explicit assumptions about diurnal shape, seasonal variation, commercial demand share, and tourism-sensitive peaks.

\section*{S3. Outage-setting rationale and robustness note}
The base resilience setting assumes 4 outage events per year, each with 48 h duration, a critical-load fraction of 0.55, and diesel unavailability during the outage windows. This should be interpreted as a severe island contingency stress case that represents fuel-logistics disruption rather than a universal assumption for all outages.

The base resilience results are:
\begin{itemize}[leftmargin=1.5em]
  \item battery\_only: CLSR = 1.000, EENS = 0.0 MWh, survivable outage duration = 48 h
  \item diesel\_battery: CLSR = 0.643, EENS = 45.7 MWh, survivable outage duration = 3 h
  \item battery\_hydrogen: CLSR = 1.000, EENS = 0.0 MWh, survivable outage duration = 48 h
\end{itemize}

Reviewer-facing robustness message: the diesel-unavailable setting is used to test whether hydrogen materially changes the resilience envelope under severe local-resource constraints. The paper's claim is therefore conditional and screening-oriented. It does not assert that diesel is unavailable in every real outage, nor that the reported ranking would be identical under all diesel-support assumptions.

\section*{S4. Screened parameter ranges and rationale}
Table~S4 summarizes the exogenous dimensions used for threshold screening.

\begin{tabularx}{\textwidth}{p{3.3cm}p{3cm}X}
\toprule
Dimension & Values used & Rationale \\
\midrule
Delivered diesel cost & 180, 380, 700 USD/MWh & Screens low, base, and high fuel-logistics conditions for remote island supply chains. \\
Outage duration & 12, 48, 96 h & Captures short, severe, and prolonged outage conditions relevant to resilience screening. \\
Hydrogen CAPEX multiplier & 0.45, 0.70, 1.00 & Tests how threshold movement responds to lower hydrogen subsystem cost assumptions. \\
Carbon price & 0, 150, 300 USD/tCO$_2$ & Represents a neutral case plus progressively stronger carbon-exposure stress. \\
\bottomrule
\end{tabularx}

At the base outage duration of 48 h and carbon price of 150 USD/tCO$_2$, the annual cost difference between battery--hydrogen and diesel--battery is positive at low diesel cost and high hydrogen CAPEX, but becomes negative when delivered diesel cost reaches 700 USD/MWh and the hydrogen CAPEX multiplier is reduced to 0.70 or 0.45. This is the threshold logic visualized in the main-text heat map.

\section*{S5. Reproducibility workflow and repository-ready contents}
The submission package is backed by a repository-ready workflow containing:
\begin{itemize}[leftmargin=1.5em]
  \item \texttt{configs/scenarios.json} for case, cost, resilience, and sensitivity assumptions
  \item \texttt{data/philippines\_offgrid\_8760.csv} and \texttt{data/synthetic\_island\_8760.csv}
  \item dispatch model modules in \texttt{model/}
  \item reproducible scripts in \texttt{scripts/}
  \item output tables in \texttt{results/}
  \item generated figures in \texttt{figures/}
\end{itemize}

The exact Windows commands documented in the repository README are:
\begin{verbatim}
.\.venv\Scripts\python scripts\generate_synthetic_data.py
.\.venv\Scripts\python scripts\run_scenarios.py
.\.venv\Scripts\python scripts\run_sensitivity.py
.\.venv\Scripts\python scripts\plot_results.py
\end{verbatim}

This supplementary file is intended to reduce reviewer ambiguity about portfolio definitions, scenario inputs, and the reproducibility status of the paper. A public repository release can be created upon acceptance without changing the scientific content of the submitted manuscript.

\end{document}
